
\subsubsection{Problem Formulation}

Standard DPO learns from preference pairs (x, y\_w, y\_l) where y\_w is preferred over y\_l for prompt x. The hypothesis is that response text contains implicit agency signals---patterns indicating how the response engages users---that are orthogonal to preference labels.

BiDPO augments DPO with an auxiliary agency objective:

$$\mathcal{L}_{\text{BiDPO}} = \mathcal{L}_{\text{DPO}} + \lambda \cdot \mathcal{L}_{\text{agency}}$$

where $\lambda \in$ [0, 1] controls the relative weight of agency preservation.

\subsubsection{Collaboration Score}

A heuristic collaboration score quantifies agency-preserving patterns in response text:

$$\text{collab\_score}(y) = \frac{\text{raw\_score}(y)}{\sqrt{|y|_{\text{words}}}}$$

The square-root normalization by word count prevents bias toward verbose responses. The raw score aggregates four pattern categories detected via regular expressions:

\textbf{Reasoning Traces:} Explicit causal connectives (''because,'' ''therefore,'' ''since,'' ''this means,'' ''as a result,'' ''so that'').

\textbf{Uncertainty Acknowledgment:} Epistemic markers (''I think,'' ''might,'' ''could be,'' ''perhaps,'' ''uncertain,'' ''possibly'').

\textbf{Engagement Cues:} User-directed phrases (''you could,'' ''consider,'' ''option,'' ''you might,'' ''what do you'') and question marks.

\textbf{Explanation Depth:} Enumeration markers (''first,'' ''second,'' ''third,'' ''step N,'' numbered lists, bullet points).

\subsubsection{Agency Loss}

The agency loss is computed as:

$$\mathcal{L}_{\text{agency}} = (1 - \text{collab\_norm}(y_w)) - 0.5 \cdot (1 - \text{collab\_norm}(y_l))$$

where collab\_norm clips and normalizes the raw score to [0, 1]. This formulation penalizes low collaboration scores in chosen responses while providing a smaller penalty reduction for rejected responses.

\subsubsection{Training Configuration}

Training uses Mistral-7B-Instruct-v0.2 with the HH-RLHF dataset (helpful-base split) \citep{bai2022training}. Hyperparameters: $\beta$ = 0.1, $\lambda$ = 0.5, learning rate $5\times 10^{-7}$ with cosine schedule, effective batch size 16 (gradient accumulation), 250 steps at proof-of-concept scale, bfloat16 precision, gradient clipping at 1.0.

